% large-field-kp.tex — Chapter on the large-field Kotecky-Preiss
% certificate (rho_lf, K_lf) for SU(2), derived from a constructive
% character-tail truncation of the heat-kernel partition function.
%
% Linked code / Lean:
%   - computations/large_field_kp.py
%   - lean/YMGap/Sheaf/SU2KP.lean (K_lf canonical / Wilson derivations)

\chapter{The Large-Field Kotecký--Preiss Certificate}
\label{ch:large-field-kp}

\noindent\emph{The large-field activity-decay rate $\rholf$, the
cluster-expansion summability constant $\Klf$, and the constructive
character-tail derivation that closes them in $\Qp$ at each pilot
operating point.}

This chapter derives the large-field Kotecký--Preiss certificate
$(\rholf, \Klf)$ for $\SUtwo$ at scale $k$ from a constructive
truncation of the heat-kernel character expansion of
\Cref{ch:heat-kernel}, equation~\eqref{eq:heat-kernel-su2}.  The
construction is structurally independent of the small-field
machinery of \Cref{ch:small-field-kp}: the large-field certificate
needs only the per-plaquette activity bound away from the identity
on $\SUtwo$ and the plaquette-polymer connectivity constant
$\muplaq$ on the 4D lattice.  The derivation replaces a placeholder
$\Klf = 1$ used in earlier project iterations with an explicit
$\Qp$-rational closure verified at the canonical and Wilson pilot
operating points.

\section{The character-tail partition function}
\label{sec:Zt}

\marginkey{The non-trivial-irrep tail $Z_t$}%

Define
\begin{equation}
  \label{eq:zt-definition}
  Z_t \;:=\; \sum_{j \in \tfrac12\mathbb{N}_{> 0}} (2j+1)^2 \cdot e^{-t\,\Casimir(j)}
  \;=\;
  \sum_{2j \ge 1} (2j+1)^2 \cdot e^{-t\,j(j+1)} ,
\end{equation}
the partition function on the non-trivial $\SUtwo$ irreps with the
heat-kernel time parameter $t = t(g_k)$ of \Cref{ch:heat-kernel}.
The leading term is $j = \tfrac12$, contributing $4\,e^{-3t/4}$.

\begin{proposition}[Truncation with geometric tail]
\label{prop:zt-truncation}
For any cutoff $J \in \tfrac12\mathbb{N}$ and every $t \ge 1$,
\begin{equation}
  \label{eq:zt-truncation}
  Z_t
  \;=\;
  \sum_{2j = 1}^{2J} (2j+1)^2 \cdot e^{-t\,j(j+1)}
  \;+\; \mathrm{tail}_{J}(t) ,
\end{equation}
with
\begin{equation}
  \label{eq:zt-tail-bound}
  \mathrm{tail}_{J}(t)
  \;\le\;
  (2J + 3)^2 \cdot e^{-t\,(J + \tfrac12)(J + \tfrac32)}
  \cdot \frac{1}{1 - r_J(t)},
  \qquad
  r_J(t) := e^{-t\,(2J + 2)} \cdot \left(\frac{2J+5}{2J+3}\right)^2 ,
\end{equation}
the geometric tail sum bounded by the worst-case ratio between
consecutive terms.  Truncation at $J = 6$ keeps $\mathrm{tail}_6(t)
< 10^{-15}$ for every $t \ge 1$.
\end{proposition}

\begin{proof}
The tail is $\sum_{2j > 2J} (2j+1)^2 e^{-t\,j(j+1)}$, dominated by
the term $j = J + \tfrac12$ which contributes
$(2J+3)^2\,e^{-t(J+\tfrac12)(J+\tfrac32)}$.  The ratio between
consecutive terms in the tail is bounded above by $r_J(t)$, so the
remaining contribution is bounded by the geometric series
$1/(1 - r_J(t))$.  At $J = 6$ and $t \ge 1$ the prefactor is $13^2
\cdot e^{-42}$, far below $10^{-15}$.
\end{proof}

The $J = 6$ truncation is the manuscript-standard cutoff
(it is also the default in the companion Python module);
downstream estimates use the truncated $Z_t$ replaced by an
explicit $\Qp$-rational upper bound $Z_t^{\mathrm{upper}}$ that
includes the geometric tail correction.

\section{The large-field activity bound}
\label{sec:lf-activity}

\marginkey{Per-plaquette tail bound}%

\begin{theorem}[Large-field activity bound]
\label{thm:large-field-activity}
For every polymer $X \in \Pcal_k$ supported in the large-field
region $\Omklf$,
\begin{equation}
  \label{eq:large-field-activity}
  |a_k^{\mathrm{lf}}(X)| \;\le\; Z_t^{|X|},
\end{equation}
where $|X|$ is the plaquette count of $X$ and $Z_t$ is the
character-tail partition function of \eqref{eq:zt-definition}.
\end{theorem}

\begin{proof}[Proof sketch]
Each plaquette in the large-field region contributes a factor of
$Z_t$ under the heat-kernel character expansion of
\Cref{ch:heat-kernel}, equation~\eqref{eq:heat-kernel-su2},
restricted to non-trivial irreps.  The polymer-activity bound is
the product over plaquettes in $X$.  See
\cite{DrouffeItzykson1978} \S3 and \cite{ItzyksonDrouffe1989}
\S5 for the standard heat-kernel character-expansion treatment.
\end{proof}

The activity-decay rate is therefore
\begin{equation}
  \label{eq:rho-lf}
  \rholf(g_k) \;:=\; -\log Z_t^{\mathrm{upper}}(g_k),
\end{equation}
positive in the regime $Z_t^{\mathrm{upper}}(g_k) < 1$, equivalently
$t \ge t^\star$ for an explicit threshold $t^\star$.

\section{The plaquette-polymer connectivity constant
  \texorpdfstring{$\muplaq$}{mu\_plaq}}
\label{sec:mu-plaq}

The Brydges--Federbush convergence theorem applied to the
large-field polymer expansion requires a polymer-counting bound on
the plaquette graph of the 4D hypercubic lattice.

\begin{convention}[Plaquette connectivity, $\muplaq$]
\label{conv:mu-plaq}
The number of connected plaquette-polymers of size $n$ containing a
fixed plaquette is bounded by $\muplaq^n$, with
\begin{equation}
  \label{eq:mu-plaq-default}
  \muplaq \;:=\; 16
  \quad (\textnormal{literature-shape}, \cite{Seiler1982});
\end{equation}
a genuinely Brydges--Federbush-conservative fallback
$\muplaq^{\mathrm{BF}} = 20\,e \approx 54.4$ is available from the
direct combinatorial edge-degree count
\cite{BrydgesFroehlichSeiler1978IHES}.
\end{convention}

The two values are in tension: $20\,e \approx 54$ is rigorously
derivable from a direct count (each plaquette has $4$ edges, each
edge sits in $2(d - 1) - 1 = 5$ other plaquettes for $d = 4$, giving
plaquette edge-degree $20$ and BF-conservative
$\muplaq^{\mathrm{BF}} = e \cdot 20 = 20\,e$), while $5d - 4 = 16$
is the standard lattice-gauge value attributed to Seiler that
exploits planarity beyond the universal BF bound.  The manuscript
adopts $\muplaq = 16$ as the default; downstream sensitivity is
audited in \Cref{ch:appx-fiducials}.

\section{The cluster-expansion summability constant
  \texorpdfstring{$\Klf$}{K\_lf}}
\label{sec:klf}

\marginkey{Brydges--Federbush convergence}%

\begin{theorem}[Large-field cluster-expansion convergence]
\label{thm:klf-convergence}
Set $q := \muplaq \cdot Z_t^{\mathrm{upper}}$.  The large-field
polymer sum
\begin{equation}
  \label{eq:polymer-sum-lf}
  \sum_{X \ni p} |a_k^{\mathrm{lf}}(X)| \cdot e^{\rholf \cdot |X|}
  \;\le\;
  \sum_{n \ge 1} \muplaq^n \cdot Z_t^n \cdot e^{n\,\rholf}
  \;=\;
  \sum_{n \ge 1} q^n
\end{equation}
converges to a finite bound $\Klf < \infty$ if and only if $q < 1$,
in which case
\begin{equation}
  \label{eq:k-lf-formula}
  \Klf \;\le\; \frac{q}{1 - q} \;=\; \frac{\muplaq \cdot Z_t^{\mathrm{upper}}}{1 - \muplaq \cdot Z_t^{\mathrm{upper}}} .
\end{equation}
\end{theorem}

\begin{proof}
The polymer-activity bound \cref{thm:large-field-activity} gives
$|a^{\mathrm{lf}}(X)| \le Z_t^{|X|}$, and the connectivity bound
\cref{conv:mu-plaq} gives at most $\muplaq^n$ polymers of size $n$
through a fixed plaquette.  Multiplying by the
$e^{\rholf\,|X|}$ weight and using $e^{\rholf} = 1/Z_t^{\mathrm{upper}}$
collapses the $Z_t$ factors, leaving the geometric series in $q$.
\end{proof}

\begin{infobox}[main][title=Numerical impact]
At the canonical pilot $g_0 = 0.30$:
$Z_t^{\mathrm{upper}} \approx 9.7 \times 10^{-4}$,
$q = 16\,Z_t^{\mathrm{upper}} \approx 0.0156$,
$\Klf \approx 0.0156$.  At Wilson pilot $g_0 = 0.50$:
$Z_t^{\mathrm{upper}} \approx 2.4 \times 10^{-5}$,
$q \approx 3.9 \times 10^{-4}$, $\Klf \approx 3.9 \times 10^{-4}$.
The cluster-expansion convergence threshold $q < 1$ is comfortably
satisfied at both operating points.
\end{infobox}

\section{The certificate}
\label{sec:lf-certificate}

\begin{definition}[Large-field KP certificate]
\label{def:lf-certificate}
The \emph{large-field Kotecký--Preiss certificate at scale $k$} is
the triple
\begin{equation}
  \label{eq:lf-certificate}
  c_k^{\mathrm{lf}} \;:=\;
  \bigl(\rholf(g_k),\; \Klf(g_k),\; \mathrm{sat}_{\mathrm{lf}}\bigr)
  \;\in\; \mathrm{KPCertificate}(V, P_{\mathrm{lf}}),
\end{equation}
where $\mathrm{sat}_{\mathrm{lf}}$ is a proof of the Kotecký--Preiss
predicate
\[
  P_{\mathrm{lf}}(\rholf, \Klf) :\Leftrightarrow
  \rholf > 0 \;\wedge\;
  q := \muplaq\,e^{-\rholf} < 1 \;\wedge\;
  \Klf \le q/(1 - q),
\]
with the value type $V$ as in \Cref{ch:small-field-kp},
Definition~\ref{def:sf-certificate}.
\end{definition}

\section{Structural independence from the small-field certificate}
\label{sec:lf-independence}

The large-field certificate is derived independently of the
small-field certificate of \Cref{ch:small-field-kp}: the inputs are
$(g_k, t)$ and the plaquette-polymer connectivity constant
$\muplaq$, none of which involves $\clambda$ or the small-field
Gribov-patch geometry.  The two certificates meet only at the
\emph{buffer overlap} of \Cref{ch:overlap-kp}, where the
domination-gluing of \Cref{ch:sheaf-framework} combines them into
the overlap certificate $c_k^{\cap}$.

\section{Code provenance}
\label{sec:lf-codeanchor}

\begin{itemize}
  \item Python implementation:
    \codepath{computations/large\_field\_kp.py}.  Key entry points:
    \codepath{compute\_zt\_upper(g\_k, t\_choice, j\_max)} for the
    constructive $Z_t^{\mathrm{upper}}$ of \cref{prop:zt-truncation};
    \codepath{compute\_large\_field\_kp(g\_k, t\_choice, mu\_plaq,
    j\_max)} returns a \codepath{LargeFieldKP} record carrying
    $\rholf$, $\Klf$, the \codepath{convergent} flag of
    \cref{thm:klf-convergence}, and the four named fiducial values
    $\muplaq \in \{16,\; 8e,\; 20e,\; \text{custom}\}$.
    Tests in \codepath{tests/test\_large\_field\_kp.py}.
  \item Constants:
    \codepath{MU\_PLAQUETTE\_4D = 16} (default,
    Seiler-cited),
    \codepath{MU\_PLAQUETTE\_BF\_CONSERVATIVE\_4D = 20*e}
    (rigorously BF-conservative fallback),
    \codepath{J\_MAX\_DEFAULT = 12} (the manuscript-standard
    truncation $J = 6$ in half-integer notation; the code stores
    $2J = 12$).
  \item Lean formalization: \codepath{lean/YMGap/Sheaf/SU2KP.lean}
    carries the $\Qp$-rational closures
    \codepath{K\_lf\_canonical\_derivation} and
    \codepath{K\_lf\_wilson\_derivation} verifying the
    Brydges--Federbush convergence relation $q \le \Klf \cdot
    (1 - q)$ at the rounded $\Qp$ values of $Z_t^{\mathrm{upper}}$
    and $\muplaq$.
\end{itemize}
